\section*{Introduction}

The notes correspond to the bachelor 
course \emph{Galois theory} of the 
Vrije Universiteit Brussel, 
Faculty of Sciences, 
Department of Mathematics and Data Science. The course
is divided into twelve two-hour lectures. 

The material is somewhat standard. Basic texts on fields and Galois theory 
are for example \cite{MR1645586} and 
\cite{MR3379917}. 

As usual, we also mention a set of 
\href{https://kconrad.math.uconn.edu/blurbs/}{great expository papers} by 
Keith Conrad. The notes are extremely well-written and useful  
at every stage of a mathematical career. 

% Several chapters contain optional paragraphs that give examples of 
% how to apply \href{https://oscar.computeralgebra.de/}{OSCAR Computer Algebra System}
% to concrete problems in Galois theory. 

The notes include many exercises, some with detailed solutions. Mandatory exercises have a \colorbox{green!5!white}{green background}, while optional ones
(bonus exercises) have a \colorbox{yellow!15!white}{yellow background}.
The notes also include some additional comments. While these are entirely optional, I hope they offer further insight. They are highlighted with a \colorbox{red!5!white}{pink background}.

The notes include Magma code, which we use to verify examples and offer alternative solutions to certain exercises. Magma \cite{zbMATH01077111} is a powerful software tool designed for working with algebraic structures. There is a free \href{https://magma.maths.usyd.edu.au/calc/}{online} version of Magma available.

Several of the exercise solutions were written by Silvia Properzi.

These notes are freely available at \url{https://github.com/vendramin/galois},
where the most recent version can always be found. They are also listed in the
AMS Open Math Notes collection, reference \texttt{OMN:202508.111468}; the copy
hosted there is a snapshot and may lag behind the version on GitHub.

If you want to cite these notes, you may use the following entry:

\vspace{1cm}
{\footnotesize
\begin{quote}
\begin{verbatim}
@misc{vendraminRM,
    author       = {Vendramin, Leandro},
    title        = {Galois theory (lecture notes)},
    year         = {2026},
    howpublished = {\url{https://github.com/vendramin/galois}},
    doi          = {10.5281/zenodo.8297595},
}
\end{verbatim}
\end{quote}}
\vspace{.8cm}

\subsection*{Acknowledgements}

Wouter Appelmans, Andrew Darlington, Luca Descheemaeker, 
Alejandro de la Cueva Merino, 
Wannes Malfait, Manet Michiels, Silvia Properzi, 
Lukas Simons, Ruben Wellens. 

\vfill
\noindent\rule{\textwidth}{0.4pt}\\[2pt]
{\footnotesize\indent This version was compiled on \today\ at \currenttime.
Please send comments and corrections to
\href{mailto:Leandro.Vendramin@vub.be}{Leandro.Vendramin@vub.be}.\par}



